\section{Background and Problem Setting}
\label{sec:setup}

\subsection{Prediction-powered inference and PPI++}
\label{sec:ppi-background}

Let \((X_i,Y_i)\), \(i=1,\ldots,n\), be labeled pairs and let \(\widetilde X_j\), \(j=1,\ldots,N\), be unlabeled covariates. Suppose a fixed predictor supplies \(f(X_i)\) and \(f(\widetilde X_j)\). For a mean target, the elementary prediction-powered decomposition is
\begin{align}
\widehat\mu_{\mathrm{PPI}}
&=
\frac1N\sum_{j=1}^N f(\widetilde X_j)
+
\frac1n\sum_{i=1}^n\{Y_i-f(X_i)\}.
\end{align}
When labeled and unlabeled covariates share a marginal law and the two samples are independent, the prediction term and residual correction sum to the labeled-population mean for every fixed predictor \(f\). PPI++ generalizes this idea by tuning the strength of the prediction contribution. In our construction, \(\omega\) is the PPI-style scale on a mean-zero prediction control variate. By contrast, \(\lambda\) blends trial and observational outcome regressions inside an RCT-valid score; it is not the scale on the external prediction mean.

\subsection{RCT--OBS problem setting and notation}
\label{sec:rct-obs-setting}

We observe independent samples
\begin{align}
D_R
&=\{V_i^R=(X_i^R,A_i^R,Y_i^R):i=1,\ldots,n\}\sim P_R^{\otimes n},
\\
D_O
&=\{V_j^O=(X_j^O,A_j^O,Y_j^O):j=1,\ldots,N\}\sim P_O^{\otimes N},
\end{align}
where \(X\) is a covariate vector, \(A\in\{0,1\}\) is treatment, and \(Y\) is the observed outcome. For a function \(h\), write \(\mathbb P_{R,n}h=n^{-1}\sum_i h(V_i^R)\) and \(\mathbb P_{O,N}h=N^{-1}\sum_j h(V_j^O)\).

For the exact adaptive results, these displayed samples are the final-evaluation blocks. Separate RCT and OBS blocks \(D_R^{\mathrm{tr}},D_O^{\mathrm{tr}}\) train nuisance and ratio functions, and separate blocks \(D_R^{\mathrm{tu}},D_O^{\mathrm{tu}}\) select coefficients. Define the pre-evaluation sigma-field
\begin{align}
\mathcal F
&=
\sigma\!\left(
D_R^{\mathrm{tr}},D_O^{\mathrm{tr}},
D_R^{\mathrm{tu}},D_O^{\mathrm{tu}},
\text{prespecified split randomness}
\right).
\label{eq:pre-evaluation-field}
\end{align}
The split assignments are fixed in advance or independently randomized, with their randomness included in \(\mathcal F\). Conditional on \(\mathcal F\), every fitted nuisance, ratio, basis, and selected coefficient is fixed and
\begin{align}
\mathcal L(D_R,D_O\mid\mathcal F)
&=
P_R^{\otimes n}\otimes P_O^{\otimes N}.
\label{eq:evaluation-product-law}
\end{align}
For fixed-coefficient statements the tuning blocks may be empty.

Let \(Y(1)\) and \(Y(0)\) denote potential outcomes. The RCT-population conditional and average treatment effects are
\begin{align}
\tau(x)
&=\mathbb E_R\{Y(1)-Y(0)\mid X=x\},
&
\theta
&=\mathbb E_R\{\tau(X)\}.
\end{align}

\begin{assumption}[Trial identification]
\label{ass:trial}
Consistency holds, \(\{Y(0),Y(1)\}\perp A\mid X\) under \(P_R\), and the known trial propensity \(e(x)=P_R(A=1\mid X=x)\) satisfies \(\epsilon<e(X)<1-\epsilon\) almost surely for a constant \(\epsilon>0\).
\end{assumption}

\begin{assumption}[Common covariate marginal]
\label{ass:common-marginal}
\label{ass:covariate}
The trial and observational covariates have the same marginal law, \(P_R^X=P_O^X\). This assumption does not equate the treatment or outcome mechanisms across sources, and the observational treatment contrast may remain confounded.
\end{assumption}

\begin{assumption}[Exact covariate transport]
\label{ass:exact-transport}
Under covariate shift, \(P_R^X\ll P_O^X\) and the exact density ratio
\begin{align}
r(x)
&=\frac{dP_R^X}{dP_O^X}(x),
&
\mathbb E_O\{r(X)h(X)\}
&=\mathbb E_R\{h(X)\}
\end{align}
is available for every required integrable function \(h\). The common-marginal model is the special case \(r=1\).
\end{assumption}

\begin{assumption}[Honest evaluation and moments]
\label{ass:honest-evaluation}
Every fitted object used in an exact identity is \(\mathcal F\)-measurable, the final observations satisfy Equation~\eqref{eq:evaluation-product-law}, and every displayed conditional moment is finite almost surely. Square-integrability and stronger moment conditions are stated locally when variance estimation or Wald inference requires them.
\end{assumption}

The exact-r results use Assumption~\ref{ass:exact-transport}; their common-marginal corollaries use Assumption~\ref{ass:common-marginal} and set \(r=1\). Neither version assumes that the OBS treatment contrast is causal.

Let \(\widehat\mu_R=(\widehat\mu_{R,0},\widehat\mu_{R,1})\) and \(\widehat\mu_O=(\widehat\mu_{O,0},\widehat\mu_{O,1})\) be outcome regressions fitted on the two sources. Define
\begin{align}
g(x)
&=\widehat\mu_{O,1}(x)-\widehat\mu_{O,0}(x),
\\
\varphi(V;q)
&=q_1(X)-q_0(X)
+\frac{A}{e(X)}\{Y-q_1(X)\}
-\frac{1-A}{1-e(X)}\{Y-q_0(X)\}.
\end{align}
\begin{proposition}[RCT-validity for an arbitrary outcome regression]
\label{prop:pseudo-validity-active}
\label{prop:pseudo-validity}
Under Assumptions~\ref{ass:trial} and \ref{ass:honest-evaluation}, every fixed \(q\) with finite first moments satisfies
\begin{align}
\mathbb E_R\{\varphi(V;q)\mid X\}
&=\tau(X).
\label{eq:pseudo-valid}
\end{align}
\end{proposition}
For \(\lambda\in\mathbb R\), set
\begin{align}
\widehat\mu_\lambda
&=(1-\lambda)\widehat\mu_R+\lambda\widehat\mu_O,
&
Z_\lambda
&=\varphi(V;\widehat\mu_\lambda)
=Z_0+\lambda\Delta,
&
\Delta
&=Z_1-Z_0.
\end{align}
Equation~\eqref{eq:pseudo-valid} implies \(\mathbb E_R(Z_\lambda\mid X)=\tau(X)\) and \(\mathbb E_R(\Delta\mid X)=0\). The observational contrast \(g\) is used only as a prediction and need not equal \(\tau\).

\subsection{Existing data-fusion estimators}
\label{sec:existing-fusion}

Three existing strategies provide the nearest comparisons. First, shrinkage methods combine an unbiased experimental estimate with a potentially biased observational estimate and trade variance against observational bias \citep{rosenman2020shrinkage}. Second, adaptive RCT--OBS estimators estimate source-specific conditional effects and downweight the observational estimator when its discrepancy from the trial is large \citep{cheng2021adaptive}. Third, direct multi-source learners construct weighted pseudo-outcomes for heterogeneous effects across sources \citep{li2022directfusion}. Our estimator differs at the level of the fusion object: \(\lambda\) changes the nuisance regression inside a trial-valid score, while \(\omega\) centers a fixed prediction across study covariate samples. Neither channel treats the observational treatment contrast as identified causal truth.

Within the adaptive and integrative line, the elastic procedure of \citet{yang2023elastic} tests whether real-world data are sufficiently compatible with the randomized trial before pooling information for prespecified effect modifiers. The confounding-function approach of \citet{yang2025confounding} instead introduces and jointly estimates an HTE function and an observational confounding function under explicit cross-source identification conditions. These methods use observational outcomes through a poolability decision or an identified confounding-function model; our two fusion channels retain the RCT-valid score as the causal anchor.
